% Source: https://artofproblemsolving.com/community/c3h2354818_trig_confusion
% Stable identity: aops-topic-2354818-q01
Label the right triangle as in the figure.

\begin{asy}
usepackage("amsmath");
size(4cm);
defaultpen(fontsize(10pt));
pair P=(0,0), Q=(4,0), R=(4,3);
draw(P--Q--R--cycle,linewidth(0.9));
real m=0.3; draw(Q+(-m,0)--Q+(-m,m)--Q+(0,m));
draw(arc(P,0.8,0,degrees(atan2(3,4)))); label("$\theta$",P+1.15*dir(18));
label("$a$",(2,0),S); label("$o$",(4,1.5),E); label("$h$",(2,1.5),NW);
\end{asy}

\textbf{(a)} By the Pythagorean theorem $o^2+a^2=h^2$. Then
\[\sin^2\theta+\cos^2\theta=\dfrac{o^2}{h^2}+\dfrac{a^2}{h^2}=\dfrac{o^2+a^2}{h^2}=\dfrac{h^2}{h^2}=1,\]
and
\[\tan\theta+\cot\theta=\dfrac oa+\dfrac ao=\dfrac{o^2+a^2}{ao}=\dfrac{h^2}{ao}=\dfrac ha\cdot\dfrac ho=\sec\theta\csc\theta.\]

\textbf{(b)} After the substitution, both sides of an identity become expressions in $o$, $a$, $h$. For any positive $o$ and $a$ there is a right triangle with these legs, and its hypotenuse is $h=\sqrt{o^2+a^2}$; so $o$ and $a$ can be anything, and the only relation that always holds between the three lengths is $o^2+a^2=h^2$. Hence a true identity, after clearing denominators, must follow from this one relation, and the computation keeps arriving at it. This is not a sign of an error; it is exactly the reason the identity holds.

The method uses a right triangle with angle $\theta$, which exists only for $0^\circ<\theta<90^\circ$. For other angles the trigonometric functions are defined with the unit circle, so this argument does not cover them; the same algebra works there with the point $(\cos\theta,\sin\theta)$ on the unit circle and the relation $x^2+y^2=1$ in place of $o^2+a^2=h^2$.
